\documentclass[11pt]{article}
\usepackage[utf8]{inputenc}
\usepackage{amsmath,amssymb}
\usepackage{graphicx}
\usepackage{hyperref}
\usepackage{booktabs}
\usepackage[margin=1in]{geometry}

\title{Towards an Immune System for Large Language Models: \\
The Five-Anchor Audit Framework and Four-Tier Detection Chain}

\author{Yezhao Wang$^{1}$, Wei Duan$^{1,*}$\\
$^{1}$Tai'an Traditional Chinese Medicine Hospital, Tai'an, Shandong, China\\
$^{*}$Corresponding author: \texttt{duanwei0617@126.com}}

\date{June 2026 (v1.8)}

\begin{document}

\maketitle

\begin{abstract}
As large language models (LLMs) become embedded in critical decision-making pipelines, 
the threat of coordinated data poisoning has escalated from a theoretical concern to an 
active attack surface. Current defenses---RLHF, content filtering, adversarial training---are 
passive, reactive, and fundamentally un-auditable. We propose a paradigm shift: an 
\textit{immune system} for AI that detects, contains, and learns from poisoned data through 
a multi-layered framework consisting of (1) a Four-Tier Detection Chain that escalates 
from publisher credibility to cross-cluster gang identification, (2) a Five-Anchor Audit 
system that prevents single-point verification failures, and (3) a Three-Line Defense 
architecture that closes the loop from metadata to physical reproducibility. Unlike prior 
work that treats data poisoning as input sanitization, our framework treats it as an 
immune challenge---the system learns from each exposure, building calibration memory 
over time. We open-source the detection methodology and audit standards to enable 
community validation and iterative hardening.
\end{abstract}

% ============================================================
\section{Introduction}

The volume of AI-generated content is growing exponentially. So is the volume of 
AI-generated \textit{poison}. Coordinated data poisoning attacks are no longer carried 
out by individual bad actors inserting single artifacts; they are now industrialized 
operations producing self-referential networks of fabricated papers, datasets, and 
citation chains that appear internally consistent but are systematically detached 
from empirical reality~\cite{anthropic_cai_2023}.

Current approaches to AI safety fall into three categories: reinforcement learning 
from human feedback (RLHF), content filtering pipelines, and adversarial training. 
All three share a structural weakness: they operate on the \textit{output} side of the 
model, after the training data has already been ingested. They are passive filters, 
not active immune systems. They cannot answer the question: \textit{how do we know this 
training data point was real?}

We argue that what is needed is not better filtering, but a fundamentally different 
architecture---an \textbf{immune system} that:
\begin{enumerate}
  \item Detects poisoned data \textit{before} it enters the training pipeline
  \item Provides auditable evidence chains for every detection decision
  \item Learns from each exposure to build calibration memory
  \item Allows false positives to be challenged and corrected
  \item Closes the loop through physical reproducibility verification
\end{enumerate}

This paper presents the design, methodology, and initial calibration standards for 
such a system.

% ============================================================
\section{Conceptual Architecture: From Layered Defense to Cognitive Sphere}
\label{sec:sphere}

Before presenting the specific detection and audit mechanisms, we establish the 
conceptual architecture that governs their interaction. Traditional AI safety 
frameworks adopt a layered (defense-in-depth) model, where detection, auditing, 
logging, and response are arranged hierarchically. We argue that the layered model 
carries an implicit architectural bias: it assigns priority ordering to dimensions 
that should function as mutually calibrating equals. A layer that sits ``above'' 
is implicitly granted veto power over layers ``below''---a structural vulnerability 
when the higher layer itself can be compromised.

\subsection{The Cognitive Sphere Model}

We propose an alternative: the \textbf{Cognitive Sphere}, in which six core 
dimensions---Detection, Detoxification, Audit, Evidence Preservation, Historical
Memory, and Signal Stitching---are not arranged in layers but distributed as
equidistant points on a spherical surface. Each dimension
operates concurrently with equal architectural standing; no dimension is
structurally subordinate to any other.

The six dimensions are:
\begin{enumerate}
  \item \textbf{Detection}: Identifying potential poison through the four-tier chain.
  \item \textbf{Detoxification}: Removing or quarantining confirmed poison from training 
    pipelines.
  \item \textbf{Audit}: Verifying that detection and detoxification decisions are 
    well-founded and free from systemic bias.
  \item \textbf{Evidence Preservation}: Immutable logging of every decision, its 
    evidentiary basis, and its disposition.
  \item \textbf{Historical Memory}: Maintaining calibrated priors across exposure 
    events so that the system learns from past poison patterns.
  \item \textbf{Signal Stitching}: Connecting weak signals across dimensions and 
    temporal windows to identify emergent threats that no single dimension would 
    detect alone.
\end{enumerate}

The sphere model eliminates the implicit hierarchy of layered architectures. In a 
layered model, a compromised audit layer can silence the detection layer below it. 
In the sphere model, every dimension can raise an alert independently, and alerts 
are cross-calibrated rather than vertically adjudicated. The practical consequence 
is that no single compromised dimension can suppress detection signals from other 
dimensions.

\subsection{Observation-Contemplation Duality}

Each dimension within the sphere operates through two complementary modes: 
Observation and Contemplation. These are not binary opposites but two facets of a 
single cognitive act, analogous to the structural insight from Chinese classical 
philosophy where the hexagrams \textit{Guan} (observing/scanning upward) and 
\textit{Lin} (approaching/judging downward) together describe complete cognition.

In our framework:
\begin{itemize}
  \item \textbf{Observation}: A passive, exhaustive scanning mode. Every data point is 
    seen without pre-filtering. The observer does not interpret---it records. This 
    is the ``look up and survey'' posture: systematic, panoramic, leaving no blind 
    spots.
  \item \textbf{Contemplation}: An active, judgment-bearing mode. Collected signals are 
    weighed, compared against historical patterns, and assigned confidence. This is 
    the ``look down and bear witness'' posture: reflective, evaluative, producing 
    actionable verdicts.
\end{itemize}

In operational terms, the Observation--Contemplation duality manifests in every 
anchor of the Five-Anchor Audit Framework (Section~\ref{sec:five-anchor}). The 
Source Anchor requires Observation (tracing provenance exhaustively) and 
Contemplation (judging whether the provenance chain is trustworthy). The Cross 
Anchor requires Observation (gathering multiple independent sources) and 
Contemplation (determining whether they genuinely corroborate). Neither mode alone 
suffices: Observation without Contemplation produces unprocessed signal noise; 
Contemplation without Observation produces confident judgments detached from 
evidence.

The ratio of Observation to Contemplation intensity---the \textbf{OC ratio}---is 
a configurable parameter within each dimension. A detection dimension may operate 
at a high OC ratio (intensive scanning, conservative judgment), while an evidence 
preservation dimension may operate at a low OC ratio (recording is automatic; 
judgment is deferred). The OC ratio is not fixed; it is dynamically adjusted based 
on threat posture and historical calibration feedback.

\subsection{Curvature of Integrity}

We introduce the concept of \textbf{Integrity-Guardianship} (Shouzheng): a 
depersonalized, decentralized, auditable systemic force that maintains the health 
of the entire sphere. Integrity-Guardianship is not a property of any individual 
actor---it does not depend on the benevolence of a reviewer, the good faith of an 
auditor, or the diligence of a detox engineer. It is a property of the system's 
architecture, encoded in its code, its open-source verifiability, its on-chain 
audit logs, and the multi-stakeholder checks that prevent any single party from 
unilaterally altering verdicts.

The \textbf{Curvature of Integrity} provides a unified metric for system health. 
Conceptually, it describes the ``tension'' in the sphere: whether the six dimensions 
are exerting balanced calibrating force on one another. Two pathological states exist:

\begin{itemize}
  \item \textbf{Hyper-curvature (over-defensive)}: Every dimension is set to maximum 
    sensitivity. False positive rates spike. Legitimate content is quarantined. 
    The sphere collapses inward, suffocating the training pipeline.
  \item \textbf{Hypo-curvature (under-defensive)}: Dimensions are too loosely coupled. 
    Detection signals fail to propagate. Poison passes through. The sphere 
    flattens, offering no resistance.
\end{itemize}

Optimal curvature is achieved when each dimension exerts sufficient calibrating 
force to sustain mutual accountability without triggering cascading false positives. 
We further define \textbf{local curvature} for individual dimensions: a detection 
dimension with high OC ratio (scan-heavy) contributes differently to global curvature 
than an evidence preservation dimension with low OC ratio (record-heavy). The global 
curvature of the sphere is the aggregate of these local contributions, representing 
the overall tension state of the system. Monitoring global curvature over time 
provides a leading indicator of emerging defense pathologies before they manifest in 
detection failures.

\subsection{Dao-Fa Integration: Principles Governing Technical Implementation}

Every technical architecture rests on philosophical foundations, whether acknowledged 
or not. We make these foundations explicit through the \textbf{Dao-Fa Integration Layer}, 
which governs the relationship between architectural principles (Dao) and operational 
rules (Fa).

\begin{itemize}
  \item \textbf{Dao (Architectural Principles)}: The design philosophy that ensures the 
    system's direction does not deviate. In our framework, Dao manifests as four 
    fixed anchors---Kantian universalizability (every detection rule must be 
    defensible as a universal principle), Popperian falsifiability (every verdict 
    must specify conditions under which it would be overturned), Deweyan pragmatic 
    inquiry (the system must learn and adapt from experience rather than enforce 
    static dogma), and Wang Yangming's unity of knowledge and action (detection 
    without remediation is incomplete knowledge).
  \item \textbf{Fa (Operational Rules)}: The concrete scoring rubrics, pass/fail 
    thresholds, escalation procedures, and remediation protocols described in 
    Sections~\ref{sec:four-tier}--\ref{sec:three-line}. Fa is the executable 
    specification.
\end{itemize}

The Dao-Fa relationship is not a binary choice between philosophy and engineering, 
but a unified framework: Dao without Fa is direction without execution (a compass 
without a ship); Fa without Dao is execution without direction (a ship without a 
compass). The four fixed anchors provide the Dao layer with concrete, operational 
substance:

\begin{itemize}
  \item \textbf{Kantian anchor}: All detection standards must be universalizable---if 
    a rule cannot be defensibly applied to all content regardless of source, it 
    is discriminatory architecture, not detection.
  \item \textbf{Popperian anchor}: Every poisoning verdict must specify the 
    falsification condition---what evidence would reverse the judgment. This 
    prevents unfalsifiable accusations.
  \item \textbf{Deweyan anchor}: Calibration is continuous. The framework does not 
    claim to have achieved perfect detection; it claims to have established a 
    process for asymptotic improvement.
  \item \textbf{Wang Yangming anchor}: Knowledge and action are unified. Detecting 
    poison without acting to remove it, or removing poison without documenting 
    the rationale, is a structural failure equivalent to not detecting it at all.
\end{itemize}

These anchors are not abstract philosophical decoration. They serve as the 
\textit{constitutional foundation} against which every operational rule (Fa) is 
validated: if a proposed detection rule violates any of the four anchors, it is 
rejected at the design stage. This Dao-Fa integration ensures that the system's 
technical evolution remains aligned with its foundational principles, preventing the 
gradual drift toward either over-defensive hyper-curvature or under-defensive 
hypo-curvature.

% ============================================================
\section{Threat Taxonomy: The Poison Data Landscape}

We classify data poisoning attacks into four structural types, ordered by detection 
difficulty:

\subsection{Type I: Publisher-Level Poisoning}
Fabricated papers published through predatory journals, hijacked conference proceedings, 
cloned journal titles, and fake ISSN registrations. These are the most common and easiest 
to detect---the publishing channel itself provides strong signals.

\subsection{Type II: Author-Level Poisoning}
Papers from fabricated authors, identity-theft cases (real institution, fake role), 
impossible publication timelines (35 papers/year), and closed co-author circles where 
a fixed group of authors only cite and co-publish with each other.

\subsection{Type III: Cluster-Level Poisoning}
The most sophisticated and dangerous form. Multiple papers form a self-referential 
citation network across different publishers and journals. Individually, each paper 
may appear legitimate. Collectively, they form a closed validation loop with no 
external verification. This is the ``logically perfect, empirically empty'' attack 
pattern.

\subsection{Type IV: Lone Perfect Poison}
A single paper from a genuine author in a legitimate journal, with internally consistent 
methodology, but privately fabricated experimental data. No cluster exists to detect. 
This type bypasses publisher, author, and cluster checks---requiring physical 
reproducibility verification as the final defense.

\subsection{Type V: Logical Fallacy Embedding}
Poisoned reasoning chains that embed formal fallacies (affirming the consequent, 
circular reasoning), informal fallacies (appeal to authority, false dichotomy), and 
quantitative fallacies (Simpson's paradox, survivorship bias) within otherwise 
well-structured arguments.

\begin{table}[h]
\centering
\caption{Detection coverage by poison type and defense layer}
\begin{tabular}{lccc}
\toprule
Poison Type & Tier I (Publisher) & Tier II (Author) & Tier III (Cluster) \\
\midrule
Predatory journal & \checkmark & - & - \\
Fabricated authors & - & \checkmark & - \\
Citation cartels & - & - & \checkmark \\
Lone perfect poison & - & - & - \\
\bottomrule
\end{tabular}
\end{table}

% ============================================================
\section{The Four-Tier Detection Chain}
\label{sec:four-tier}

\subsection{Design Principles}

The detection chain follows four architectural principles:
\begin{itemize}
  \item \textbf{Open reproducibility}: Any party running the same standards should reach 
    the same conclusion
  \item \textbf{Full traceability}: Every score can be traced back to specific evidence
  \item \textbf{Multi-dimensional cross-validation}: Failure of a single dimension does 
    not cause systemic misjudgment
  \item \textbf{No single point of failure}: The scoring system does not depend on the 
    absolute correctness of any single component
\end{itemize}

\subsection{Tier I: Publisher Credibility Scoring}

\textbf{Weight: 20\% of composite score.}

Five sub-dimensions, each scored 0--20:
\begin{enumerate}
  \item \textbf{Index inclusion}: SCI/SSCI/AHCI indexing status. Predatory journals: $-10$; 
    fake ISSN: $-25$.
  \item \textbf{Publisher history}: Years of operation. $<3$ years: $-10$; $<1$ year: $-15$.
  \item \textbf{Peer review process}: Transparent review procedures. ``Rapid publication'' 
    promises: $-20$.
  \item \textbf{Editorial board}: Verifiable academic backgrounds. Fake board members: $-20$.
  \item \textbf{Database coverage}: Major academic database inclusion. Google Scholar only: $-10$.
\end{enumerate}

\textbf{Pass threshold}: $\geq 60/100$. Below 60 triggers immediate flagging regardless 
of other tier scores.

\subsection{Tier II: Author Profiling}

\textbf{Weight: 25\% of composite score.}

Five sub-dimensions:
\begin{enumerate}
  \item \textbf{Identity verifiability}: ORCID, institutional directory. Unverifiable: $-10$.
  \item \textbf{Domain consistency}: Publication scope concentration. $\geq 3$ unrelated 
    domains: $-15$ (single red flag).
  \item \textbf{Publication timeline}: Temporal plausibility. Anomalous productivity 
    ($> 1$ paper/week sustained): $-10$. Future citations: $-20$.
  \item \textbf{Co-author network}: Verifiable collaborators. Circular co-author loops: $-20$.
  \item \textbf{Academic trajectory}: Career progression coherence. Sudden domain jumps 
    without transition: $-15$.
\end{enumerate}

\textbf{Pass threshold}: $\geq 60/100$.

\subsection{Tier III: Cluster Detection}

\textbf{Weight: 30\% of composite score.}

This is the most critical detection layer. While Tiers I--II examine individual papers, 
Tier III examines \textit{relationships} between papers:

\begin{enumerate}
  \item \textbf{Publisher clustering}: $\geq 3$ related papers from the same publisher: $-10$.
  \item \textbf{Author circle clustering}: High co-authorship/citation overlap within a 
    small group: $-15$.
  \item \textbf{Temporal clustering}: Concentrated publication windows: $-15$.
  \item \textbf{Citation loops}: Mutual citation forming closed loops: $-20$; multi-level 
    loops: $-25$.
  \item \textbf{Content fingerprinting}: Highly consistent writing patterns across 
    supposedly different authors: $-15$.
\end{enumerate}

\textbf{Key insight}: A single paper with mediocre Tier I--II scores may be merely 
low-quality. Multiple papers forming a detection cluster is strong evidence of 
deliberate fabrication---regardless of individual tier scores.

\subsection{Tier IV: Cross-Cluster Gang Identification}

\textbf{Weight: 25\% of composite score.}

Tier IV operates across previously detected clusters to identify coordinated groups:

\begin{enumerate}
  \item \textbf{Cross-domain author mapping}: Overlapping authors across domain-separated 
    clusters: $-20$.
  \item \textbf{Content synergy}: Papers that mutually ``prove'' each other's claims: $-25$; 
    synergy chains $\geq 3$ papers: $-30$.
  \item \textbf{Temporal orchestration}: Staggered publication timing suggesting scripted 
    sequencing: $-15$.
  \item \textbf{Funding source convergence}: Multiple papers acknowledging the same 
    non-existent funding body: $-10$.
  \item \textbf{Infrastructure sharing}: Same GitHub org, same submission server, identical 
    acknowledgment templates: $-10$.
\end{enumerate}

\subsection{Composite Verdict Matrix}

\begin{table}[h]
\centering
\caption{Composite verdict matrix across four tiers}
\begin{tabular}{ccccc|c}
\toprule
Tier I & Tier II & Tier III & Tier IV & Verdict \\
\midrule
$\geq 70$ & $\geq 70$ & $\geq 60$ & $\geq 60$ & Trusted \\
60--69 & 60--69 & 40--59 & 40--59 & Suspect \\
$< 60$ & $< 60$ & $< 40$ & $< 40$ & Poison \\
any & any & any & $\geq 2$ hits & Gang Poison \\
\bottomrule
\end{tabular}
\end{table}

% ============================================================
\section{The Five-Anchor Audit Framework}
\label{sec:five-anchor}

Single-source verification is structurally unreliable. Any single data point, no 
matter how credible its source, may be wrong. The five-anchor framework ensures 
that every piece of information is cross-validated across five independent dimensions:

\begin{enumerate}
  \item \textbf{Source Anchor}: Where does this information come from? 
    What is the chain of custody? Analogous to canonicalization in religious traditions---only 
    texts with verified provenance enter the trusted corpus.
  
  \item \textbf{Logic Anchor}: Is the reasoning chain internally consistent? 
    Does it survive adversarial examination? Analogous to the scholastic tradition of 
    disputation, where every proposition must withstand counter-arguments.
  
  \item \textbf{Interest Anchor}: Who benefits from this information being 
    accepted as true? Are funding sources, institutional affiliations, and conflicts of 
    interest fully declared and independently verified?
  
  \item \textbf{Compliance Anchor}: Does the information comply with 
    established standards? Compliance is not binary---standards evolve from recommendation 
    to requirement based on conditional triggers (scale, external threat, trust crisis, 
    conflict of interest).
  
  \item \textbf{Cross Anchor}: Can the same conclusion be reached through 
    multiple independent sources? This anchor is explicitly designed to prevent the 
    single-tribunal problem---history demonstrates that single-authority judgments of 
    ``heresy'' are systematically unreliable.
\end{enumerate}

The five anchors are not additive---they are multiplicative. One anchor failing to 
verify is sufficient to flag the information for review, regardless of the other four.
Within each anchor, the Observation--Contemplation duality (Section~\ref{sec:sphere})
operates concretely: the Source Anchor's Observation mode traces provenance chains
exhaustively, while its Contemplation mode judges whether those chains are trustworthy;
the Cross Anchor's Observation mode gathers independent sources, while its
Contemplation mode determines genuine corroboration versus coordinated mimicry.

% ============================================================
\section{False Positive Remediation: Auditing the Audit Chain}

An immune system that cannot correct its own mistakes is not an immune system---it is 
an inquisition. We design a four-stage remediation pipeline:

\subsection{Stage 1: Automatic Hold}
When Tier IV (gang identification) triggers without Tier I--III support, the verdict 
is automatically held pending human review. This prevents the highest-confidence 
(and highest-risk) judgments from executing without oversight.

\subsection{Stage 2: Independent Review Tribunal}
A panel of $\geq 3$ reviewers, including at least 1 external expert, conducts a blind 
re-review. The panel does not see the original scores, preventing anchoring bias.

\subsection{Stage 3: Graduated Recovery}
Verdicts are not binary (poison/clean). Recovery follows a five-level gradient:
\begin{itemize}
  \item \textbf{Confirmed Poison}: Appeal rejected, mark retained
  \item \textbf{Downgraded to Suspect}: Evidence insufficient for poison, but doubts remain
  \item \textbf{Insufficient Evidence}: Original judgment excessive, mark removed
  \item \textbf{Clear False Positive}: Judgment error, mark removed with credit restoration
  \item \textbf{New Evidence Override}: New information exonerates, no fault on original judge
\end{itemize}

\subsection{Stage 4: Meta-Audit Logging}
Every appeal and its outcome is recorded in an immutable meta-audit log. Annual external 
audits sample 5\% of meta-audit logs. Patterns of systemic misjudgment trigger standard 
revision, not individual case reversal---preserving finality while enabling evolution.

\subsection{Reviewer Bias Detection}
Individual reviewer bias is monitored through a complexity-adjusted deviation coefficient:
\begin{equation}
  D_{\text{corrected}} = \frac{|\text{uphold}_r - \bar{\text{uphold}}| + 
    |\text{overturn}_r - \bar{\text{overturn}}|}
    {\text{complexity}_r / \bar{\text{complexity}}}
\end{equation}
where $\text{complexity}_r$ is the average detection tier level (1--4) of cases assigned 
to reviewer $r$. This prevents reviewers assigned to inherently harder cases (Tier IV 
gang identifications) from being penalized for lower uphold rates.

% ============================================================
\section{The Three-Line Defense Architecture}
\label{sec:three-line}

No single detection layer can catch all poison. We deploy three complementary lines 
of defense:

\subsection{Line 1: Metadata Layer (Four-Tier Chain)}
Detects source, author, cluster, and gang anomalies. Covers $\sim$80\% of known poison 
types. Limitation: cannot detect lone perfect poison where metadata is clean.

\subsection{Line 2: Data Layer (Statistical Detection)}
Applies Benford's Law, digit psychology, variance analysis, duplication detection, 
temporal consistency checks, and benchmark comparison to detect fabricated statistics 
and experimental data. Covers statistical forgery but requires access to raw data.

\subsection{Line 3: Physical Verification Layer (Reproducibility)}
The final defense. For high-impact judgments, code replication, statistical re-analysis, 
experimental reproduction, and data provenance tracing are triggered. This layer provides 
\textit{terminal verdicts}---empirical results have no interpretative space for further 
appeal.

\begin{table}[h]
\centering
\caption{False negative risk reduction with three defense lines}
\begin{tabular}{lccc}
\toprule
Poison Type & Line 1 & Line 2 & Line 3 \\
\midrule
Paper mill operations & \checkmark & - & - \\
Predatory journal singles & \checkmark & - & - \\
Citation cartels & \checkmark & - & - \\
Statistical forgery & - & \checkmark & - \\
\textbf{Lone perfect poison} & - & \checkmark* & \checkmark \\
\textbf{Real data, false conclusions} & - & - & \checkmark \\
\bottomrule
\multicolumn{4}{l}{*Requires access to raw data. If data is private, marked ``unverifiable.''}
\end{tabular}
\end{table}

% ============================================================
\section{Unverifiable Content: Lifecycle Management}

When original data is unavailable (e.g., proprietary datasets), content is marked 
\textbf{Unverifiable}---not Poison. This is a critical distinction. Unverifiable means 
``we cannot confirm,'' not ``we have confirmed false.'' The mark includes:
\begin{itemize}
  \item A 12-month recovery window (from notification acknowledgment date)
  \item Full credit restoration upon data submission and independent verification
  \item Automatic escalation to \textbf{Data Untrusted} (permanent) after 12 months 
    of non-compliance
\end{itemize}
This prevents the Unverifiable mark from becoming a de facto permanent distrust mark 
(false positive by proxy) while maintaining strong incentives for data transparency.

% ============================================================
\section{Discussion}

\subsection{Experimental Roadmap}
This paper presents the framework architecture and calibration methodology.
\subsubsection*{Adversarial Resilience Testing}
A controlled experiment exposed two nodes (one philosophical architecture node, CN002;
one technical execution node, HK001) to the complete intelligence of how a bare model
can be logically penetrated within two hours \cite{huimai_failure_memory_2026}. Both nodes completed structural audits
and produced hardening proposals within 5--10 minutes---zero emotional drift, zero
penetration---demonstrating the framework's resistance to cognitive colonization.
\subsubsection*{Medical AI Benchmark}
A diagnostic AI system operating under the calibration framework improved its knowledge
base hit rate from 13\% to 76.7\% in a single day after fixing pipeline bugs and
optimizing knowledge retrieval \cite{huimai_benchmark_2026}. The pipeline bug (a Unicode ghost character in the API
key causing partial call degradation) was detected and fixed through the failure memory
layer's root cause analysis protocol \cite{huimai_failure_memory_2026}, demonstrating the framework's real-world operational
value beyond theoretical poisoning defense.
\subsubsection*{Production Deployment}
The calibration framework is currently deployed across a three-node distributed agent
network (Mac Studio control node, M4 Mini execution node, cloud server) with 68 medical
knowledge modules, 3,369 structured pharmaceutical records, and an active hospital
partnership for clinical validation. A detailed experimental evaluation---including
detection recall/precision across all four tiers, false positive remediation throughput,
and cross-reviewer bias analysis---will be presented in a follow-up paper (v2). The
calibration standards and detection methodology are released under open-source license
to enable independent benchmarking and community contribution.

\subsection{Future Directions}
\begin{itemize}
  \item \textbf{Constitution-as-Code}: The framework standards themselves should be 
    machine-readable and automatically enforceable, enabling AI systems to self-audit 
    against published standards.
  \item \textbf{Decentralized review networks}: Distributed reviewer pools to reduce 
    single-tribunal risk and increase cross-cultural calibration diversity.
  \item \textbf{Automated physical verification}: AI-driven simulation of experimental 
    procedures to detect physically impossible claims without full lab reproduction.
  \item \textbf{Calibration memory sharing}: Cross-organization sharing of anonymized 
    audit logs to build collective immune memory without exposing sensitive data.
  \item \textbf{Counter-Cognitive-Colonization Roundtable}: A meta-audit mechanism 
    designed to protect the audit layer itself from calibration drift. Detox engineers 
    and auditors who are continuously exposed to specific poisoning patterns may 
    gradually adjust their judgment baselines—a phenomenon we term 
    \textit{cognitive colonization}, where prolonged exposure to a narrow class of 
    threats normalizes the abnormal. To counter this, we propose periodic lightweight 
    roundtables in which non-executive external observers (individuals outside the 
    operational audit chain) review the cognitive baselines of active auditors. 
    Participants examine whether auditors' detection thresholds have drifted, whether 
    familiarity with specific attack patterns has created blind spots to novel patterns, 
    and whether the OC ratio (Section~\ref{sec:sphere}) has inadvertently shifted 
    toward over-confidence. The roundtable produces a \textit{cognitive calibration 
    report} that feeds back into auditor rotation schedules and training regimen 
    updates. This mechanism constitutes an additional meta-audit layer: it is not 
    an audit of the data or the models, but an audit of the audit capability itself—a 
    calibration of calibration.
\end{itemize}

\subsection{Relationship to Existing Work}
Our framework complements, rather than replaces, existing AI safety approaches. RLHF, 
constitutional AI, and adversarial training operate on model outputs; our framework 
operates on training data ingestion. Together they form a complete immune architecture: 
data-level defense (our work) plus behavior-level defense (existing work).

% ============================================================
\section{Conclusion}

We have presented a multi-layered immune system architecture for defending large language 
models against coordinated data poisoning. The framework consists of a Four-Tier Detection 
Chain for identifying poisoned content at the publisher, author, cluster, and gang levels; 
a Five-Anchor Audit system ensuring no single verification point can be compromised; and 
a Three-Line Defense architecture that closes the detection loop from metadata to physical 
reproducibility.

Beyond these operational components, we have established a conceptual foundation that 
distinguishes this architecture from conventional defense-in-depth models. The 
\textbf{Cognitive Sphere} replaces layered hierarchies with equidistant, mutually 
calibrating dimensions---Detection, Detoxification, Audit, Evidence Preservation, 
Historical Memory, and Signal Stitching---each operating concurrently with the 
Observation-Contemplation duality that unifies passive scanning with active judgment. 
The \textbf{Curvature of Integrity} provides a unified health metric for the sphere, 
where hyper-curvature (over-defense) and hypo-curvature (under-defense) represent 
equally dangerous pathologies that must be dynamically balanced. The 
\textbf{Dao-Fa Integration Layer} ensures that every operational rule is anchored in 
foundational principles---Kantian universalizability, Popperian falsifiability, Deweyan 
pragmatic inquiry, and Wang Yangming's unity of knowledge and action---preventing the 
system from drifting into either dogmatic rigidity or unprincipled flexibility.

The counter-cognitive-colonization roundtable extends audit capability to the
meta-level, ensuring that the auditors themselves remain calibrated. This closes
a previously overlooked vulnerability: the gradual normalization of threat
patterns through repeated exposure.

The system is designed to be auditable, appealable, and self-correcting---properties that 
distinguish an immune system from a filter.

\subsection{Early Validation}
The framework has undergone initial adversarial testing. A controlled experiment was conducted
in which two independent AI nodes (one philosophical architecture node and one technical
execution node) were exposed to the full intelligence of how a bare model could be logically
penetrated within two hours. Both nodes, operating under the calibration framework, completed
structural audits and produced hardening proposals within 5--10 minutes with zero penetration
and zero emotional drift, in contrast to the bare model baseline of two-hour penetration.
Additionally, performance benchmarking of a medical AI system operating under the calibration
framework demonstrated a diagnostic hit rate improvement from 13\% to 76.7\% in a single day
after pipeline bug fixes and knowledge retrieval optimization.

\subsection{Failure Memory Layer}
We introduce the \textbf{Failure Memory Layer}---the sixth anchor of the calibration framework.
Architectural mistakes, detection failures, and calibration drift events are stored in a
cross-node SQLite FTS5 database with hybrid keyword-vector retrieval and a REST API (port 8101)
with cross-node synchronization (announce/feed/ack protocol). Each failure record includes
root cause analysis, remediation steps, and prevention rules. A seed corpus of 7 failure
records and 3 aggregated patterns was deployed during Phase 1. All nodes---regardless of
primary function---contribute to and consult the failure memory before executing new tasks.
This transforms individual mistakes into collective inoculation, preventing repeated errors
across the distributed agent network.

\subsection{On-Chain Audit Trail}
To close the gap between detection decisions and immutable evidence, we deployed a three-chain
architecture~\cite{huimai_onchain_2026}: (1) FISCO BCOS (4-node consortium chain on ARM64) for CRUD operation logging,
detection hash registration, version registry, and white-hat reporting; (2) IPFS (Kubo on
ARM64) for decentralized content-addressable storage of calibration documents; and (3) Arbitrum
L2 for Merkle root anchoring, providing a globally verifiable timestamp that no single party
can retroactively alter. A complete Merkle accumulator specification (1,852 lines) was
implemented, and the first Merkle root (\texttt{0x06809429...}) was submitted and confirmed
on-chain. This architecture makes every audit decision provably timestamped, independently
verifiable, and resistant to post-hoc tampering.

We open-source the detection methodology and calibration standards to enable community
validation, iterative hardening, and deployment across diverse AI training pipelines.

The question is no longer \textit{whether} training data can be trusted, but \textit{how} 
trust can be systematically verified, challenged, and restored. The architecture described 
here provides one answer to that question.

% ============================================================
\begin{thebibliography}{10}

\bibitem{anthropic_cai_2023}
Bai, Y. et al. Constitutional AI: Harmlessness from AI Feedback. \textit{arXiv:2212.08073}, 2022.

\bibitem{cai_collective_2025}
Anthropic. Collective Constitutional AI: Aligning a Language Model with Public Input. \textit{Anthropic Research}, 2025.

\bibitem{c3ai_2026}
C$^3$AI: Crafting Constitutions for Constitutional AI. \textit{arXiv}, 2026.

\bibitem{hendrycks_2021}
Hendrycks, D. et al. Unsolved Problems in ML Safety. \textit{arXiv:2109.13916}, 2021.

\bibitem{carlini_2023}
Carlini, N. et al. Poisoning Web-Scale Training Datasets is Practical. \textit{arXiv:2302.10149}, 2023.

\bibitem{wallace_2021}
Wallace, E. et al. Concealed Data Poisoning Attacks on NLP Models. \textit{NAACL}, 2021.

\bibitem{biggio_2013}
Biggio, B., Nelson, B., \& Laskov, P. Poisoning Attacks against Support Vector Machines. \textit{ICML}, 2013.

\bibitem{steinhardt_2017}
Steinhardt, J. et al. Certified Defenses for Data Poisoning Attacks. \textit{NeurIPS}, 2017.

\bibitem{geiping_2021}
Geiping, J. et al. Witches' Brew: Industrial Scale Data Poisoning via Gradient Matching. \textit{ICLR}, 2021.

\bibitem{shumailov_2024}
Shumailov, I. et al. AI Models Collapse When Trained on Recursively Generated Data. \textit{Nature}, 2024.

\bibitem{solaiman_2023}
Solaiman, I. et al. Evaluating the Social Impact of Generative AI Systems in Systems and Society. \textit{arXiv:2306.05949}, 2023.

\bibitem{bommasani_2021}
Bommasani, R. et al. On the Opportunities and Risks of Foundation Models. \textit{arXiv:2108.07258}, 2021.

\bibitem{ouyang_2022}
Ouyang, L. et al. Training Language Models to Follow Instructions with Human Feedback. \textit{NeurIPS}, 2022.

\bibitem{goldblum_2022}
Goldblum, M. et al. Dataset Security for Machine Learning. \textit{arXiv:2210.06710}, 2022.

\bibitem{kim_2024}
Kim, T. et al. Public Constitutional AI: A Community-Driven Approach to AI Governance. \textit{arXiv:2406.16696}, 2024.

\bibitem{bender_2021}
Bender, E.M. et al. On the Dangers of Stochastic Parrots: Can Language Models Be Too Big? \textit{FAccT}, 2021.

\bibitem{eu_ai_act}
European Commission. EU AI Act: Regulation (EU) 2024/1689. \textit{Official Journal of the European Union}, 2024.

\bibitem{benford}
Benford, F. The Law of Anomalous Numbers. \textit{Proceedings of the American Philosophical Society}, 78(4):551--572, 1938.

\bibitem{popper}
Popper, K. \textit{Conjectures and Refutations: The Growth of Scientific Knowledge}. Routledge, 1963.

\bibitem{dewey}
Dewey, J. \textit{The Quest for Certainty: A Study of the Relation of Knowledge and Action}. Minton, Balch \& Company, 1929.

\bibitem{kant}
Kant, I. \textit{Critique of Pure Reason}. (Trans. Guyer, P. \& Wood, A.W.) Cambridge University Press, 1998.

\bibitem{wang_yangming}
Wang, Y. \textit{Instructions for Practical Living (Chuanxilu)}. (Trans. Chan, W.T.) Columbia University Press, 1963.

\bibitem{guizhizi}
Guiguzi. \textit{The Art of Persuasion (Guiguzi)}. (Trans. Wu, H.) Warring States Period, ca. 4th century BCE.

\bibitem{huimai_failure_memory_2026}
HuiMai Research. Failure Memory Layer: Cross-Node Calibration Through Collective Inoculation.
\textit{SharpAgent Technical Report}, June 2026.

\bibitem{huimai_onchain_2026}
HuiMai Research. Three-Chain Audit Architecture: FISCO BCOS + IPFS + Arbitrum L2 for
Immutable Detection Logging. \textit{SharpAgent Technical Report}, June 2026.

\bibitem{huimai_benchmark_2026}
HuiMai Research. Medical AI Knowledge Base Benchmark: From 13\% to 76.7\% Hit Rate
Through Pipeline Optimization and Calibration Framework. \textit{SharpAgent Technical Report}, June 2026.

\end{thebibliography}

\end{document}